% Related Work - UAI version (condensed)

\paragraph{Probabilistic KG Embeddings.}
UKGE~\citep{chen2019embedding} associates confidence scores with triples; BEUrRE~\citep{chen2021probabilistic} uses box embeddings where volume indicates uncertainty; GP-KGE~\citep{choudhary2021perm} (a variational Gaussian model, not a kernel-based GP; see \S\ref{sec:background}) learns Gaussian distributions over entities.
Earlier, KG2E~\citep{he2015learning} assigns a learned Gaussian to each entity and each relation type---fixed learned parameters per relation type, not conditioned on whether a specific entity has been observed with that relation; these learned covariances therefore remain entity-level statistics that do not explicitly track entity-relation co-occurrence.
All these methods share the limitation we identify: learned variances are relation-agnostic.
Despite being trained on data containing entity-relation co-occurrence patterns, these methods do not learn relation-specific uncertainty from training alone, necessitating explicit decomposition.
Unlike open-world KGC~\citep{shi2018open}, which targets entities entirely absent from training via external text signals, novel-context OOD targets familiar entities appearing in unseen relations with no external signals---a purely structural phenomenon addressable by coverage tracking.
While relation-aware scoring functions exist (R-GCN~\citep{schlichtkrull2018modeling}, CompGCN~\citep{vashishth2020composition}), these condition \emph{predictions} on relations but still produce relation-agnostic uncertainty estimates; our decomposition addresses uncertainty independently of the scoring architecture.
Recent work has also addressed distribution shift for entity embeddings in the inductive setting~\citep{liu2025uncertainty}, where new entities absent from training lack coverage for all relations. This approach focuses on inductive generalization for new entities but does not study relation-specific uncertainty or coverage-based detection; our work is complementary in the transductive setting and provides the theoretical basis for why entity-level variance alone is insufficient. The broader landscape of KG uncertainty is surveyed in~\citet{jarnac2025uncertainty}.

\paragraph{OOD Detection.}
MC Dropout~\citep{gal2016dropout}, Deep Ensembles~\citep{lakshminarayanan2017simple}, and Energy-based methods~\citep{liu2020energy} capture model uncertainty but not structural observation patterns.
Spectral-normalized Neural Gaussian Process (SNGP)~\citep{liu2020simple} estimates distance to training data via a GP output layer; when applied using entity embeddings as inputs (as done here, producing per-entity uncertainty), it remains relation-agnostic.
Classification-oriented evidential approaches (e.g., \citep{sensoy2018evidential}) are not directly comparable to KG-tail OOD detection in this setup; we therefore evaluate MC Dropout, Deep Ensembles, Energy, and SNGP as baselines.
Graph Posterior Networks~\citep{stadler2021graph} propagate uncertainty via the graph structure but do not model relation-specific edge types, and target node classification rather than triple-level OOD detection.
Classical feature-space methods (Mahalanobis distance~\citep{lee2018simple}, energy scores) assume continuous feature geometry; KGs' discrete combinatorial structure motivates our coverage-based alternative.
In graph-structured data more broadly, GNNSafe~\citep{wu2023energy} applies energy-based detection to GNNs for OOD \emph{node} detection in semi-supervised node classification---a distinct setting from triple-level KG OOD---and \citet{wang2024bridging} provide a unified graph-theoretic framework jointly addressing OOD detection and generalization. Our work differs by targeting the knowledge graph triple classification setting, where the relational structure creates a distinct challenge: entity-level signals conflate relation-specific contexts.

\paragraph{Temporal Knowledge Graphs.}
Temporal KG methods extend embedding models with time-aware scoring functions for link prediction over evolving graphs~\citep{garcia2018learning,lacroix2020tensor}. These methods model \emph{when} facts hold but do not address uncertainty quantification or OOD detection---the problem we target. Our ground-truth temporal evaluation uses ICEWS14 and ICEWS18, which are included in the TGB 2.0 benchmark suite for temporal KGs~\citep{gastinger2024tgb2}. Our simulated temporal splits test whether methods can flag unseen entity-relation patterns, complementing temporal KG research.
Recent conformal prediction approaches provide coverage guarantees for KG predictions~\citep{zhu2025conformalized,zhu2025certainty}; concurrently, \citet{zhu2025condkgcp} propose predicate-conditional conformalized answer sets providing per-relation coverage guarantees. These approaches target \emph{prediction set} size guarantees rather than OOD \emph{detection}; our decomposition complements them by identifying which uncertainty type---semantic or structural---drives each prediction's unreliability.
\citet{ni2025trustworthy} study uncertainty-aware KG reasoning from a trustworthy AI perspective, using conformal prediction to provide coverage guarantees for KG reasoning tasks. Our work differs by targeting OOD \emph{detection} rather than coverage guarantees, and by providing a decomposition theorem identifying why entity-level methods fail on relational novelty.

\paragraph{Inductive and Foundation KGE.}
Foundation models for KG reasoning such as ULTRA~\citep{galkin2024ultra} enable zero-shot link prediction on unseen knowledge graphs with arbitrary entity and relation vocabularies. In inductive settings where new entities appear at test time, our coverage-based structural signal requires adaptation---an entity with no training triples has zero coverage for all relations. Extending CAGP's decomposition to inductive KGE is a natural direction for future work. The OOD landscape for graph-structured data more broadly is surveyed in~\citet{li2022oodgraph}; a dedicated graph OOD \emph{detection} survey~\citep{cai2025ooddetection} and a GNN uncertainty survey~\citep{wang2024uncertainty} provide further context for the specific challenges of the KG triple classification setting addressed here.

\paragraph{Deployed KG Systems.}
Industrial knowledge graphs track triple-level provenance and source confidence but do not explicitly track entity-relation coverage for uncertainty estimation~\citep{hogan2021knowledge}. YAGO assigns per-triple confidence derived from extraction confidence, not from observation patterns~\citep{suchanek2007yago}. Enterprise graph platforms provide quality validation and lineage tracking~\citep{noy2019industry} but lack mechanisms to flag novel relational contexts. Our coverage-based decomposition is thus complementary: rather than replacing provenance-based confidence, it identifies a failure mode---novel relational contexts---that provenance tracking does not address.

\paragraph{Positioning.}
Prior work treats uncertainty as a single quantity. We show KG uncertainty decomposes into semantic (entity-level) and structural (relation-specific) components, with empirically validated predictions on when each is necessary (Empirical Regularity~1).
Structural uncertainty (coverage) captures a form of epistemic uncertainty about relational context: has the model observed this particular entity-relation combination? This is distinct from the semantic epistemic uncertainty captured by $U_{\text{sem}}$~\citep{kendall2017uncertainties}.
Recent work has advanced KG uncertainty from complementary angles, including PAC-Bayesian generalization bounds for KGE~\citep{lee2024pacbayesian}. These address generalization guarantees, while our work targets the orthogonal problem of \emph{decomposing} uncertainty into relation-aware components for OOD detection.

\paragraph{Impossibility Results in Machine Learning.}
Our work contributes to a tradition of negative results that clarify fundamental limitations. The No Free Lunch theorems~\citep{wolpert1997nfl} establish that no learner dominates across all distributions; impossibility results in differential privacy~\citep{dwork2014algorithmic} and fairness~\citep{kleinberg2017inherent} reveal fundamental trade-offs between desirable properties. In uncertainty quantification, \citet{ovadia2019can} show that standard neural network uncertainty methods degrade under distribution shift---an empirical finding our theorem formalizes for the KG setting. Theorem~\ref{thm:impossibility} is an impossibility result of this flavor: it proves that \emph{any} entity-level uncertainty estimator is fundamentally blind to novel relational contexts, regardless of architecture or training procedure. The constructive implication---track entity-relation coverage---follows directly from the theorem.
